\chapter{แหล่งกำเนิดพลาสมาและการวินิจฉัยทางสเปกโทรสโกปี}
\label{chap:ch03}

% =========================================================================
% แผนบริหารการสอนประจำบทที่ 3 (Modular Lesson Plan Boxes - 2-Page Standard)
% =========================================================================

\vspace*{-6pt}
\section*{แผนบริหารการสอนประจำบทที่ 3}
\addcontentsline{toc}{section}{แผนบริหารการสอนประจำบทที่ 3}

% --- Page 1 of Lesson Plan ---
\begin{planbox}{หัวข้อเนื้อหาประจำบท}
\begin{enumerate}
    \item สถาปัตยกรรมแหล่งกำเนิดพลาสมาเย็นที่ความดันบรรยากาศ (APNTP Sources)
    \item การคายประจุแบบแผงกั้นไดอิเล็กทริกและหัวพ่นพลาสมาเจ็ทความดันบรรยากาศ
    \item การวินิจฉัยทางไฟฟ้า วงจรสมมูล และการวิเคราะห์กำลังไฟฟ้าด้วยแผนภาพลิสซาจู
    \item การวินิจฉัยด้วยแสง: สเปกโทรสโกปีการเปล่งแสงและการจำแนกแถบสเปกตรัมของ RONS
    \item การคำนวณอุณหภูมิการหมุนและอุณหภูมิการสั่นด้วยวิธีโบลต์ซมันน์พล็อต
    \item สถาปัตยกรรม Deep Learning (1D-CNN และ Transformer) สำหรับการแยกส่วนสเปกตรัม OES แบบเรียลไทม์
\end{enumerate}
\end{planbox}

\begin{planbox}{วัตถุประสงค์เชิงพฤติกรรม}
\begin{enumerate}
    \item จำแนกข้อแตกต่างของแหล่งกำเนิดพลาสมาความดันบรรยากาศชนิด DBD, APPJ, Corona และ Gliding Arc ได้
    \item วิเคราะห์วงจรสมมูลของระบบ DBD และคำนวณหากำลังงานไฟฟ้าที่สูญเสียจริงจากแผนภาพลิสซาจูได้อย่างถูกต้อง
    \item ระบุตำแหน่งความยาวคลื่นของการแผ่รังสีของอนุมูลอิสระ $\cdot\ce{OH}$, \ce{N2(C-B)}, และอะตอม \ce{O} จากสเปกตรัม OES ได้
    \item คำนวณอุณหภูมิของก๊าซจากอุณหภูมิการหมุนของโมเลกุลโดยใช้วิธีทางคณิตศาสตร์ Boltzmann Plot ได้อย่างแม่นยำ
    \item อธิบายหลักการทำงานของโครงข่ายประสาทเทียม 1D-CNN ในการแยกสัญญาณรบกวนและสเปกตรัมทับซ้อนในระบบ OES ได้
\end{enumerate}
\end{planbox}

\newpage

% --- Page 2 of Lesson Plan ---
\begin{planbox}{วิธีสอนและกิจกรรมการเรียนการสอน}
\begin{enumerate}
    \item บรรยายเชิงทฤษฎีควบคู่กับการสาธิตการทำงานจริงของเครื่องกำเนิดพลาสมา DBD และ APPJ ในห้องปฏิบัติการ
    \item ฝึกปฏิบัติการต่อวงจรตรวจวัดสัญญาณกระแส-แรงดันไฟฟ้าด้วยหัววัดแรงดันสูงและคาปาซิเตอร์ตรวจจับบนออสซิลโลสโคป
    \item วิเคราะห์แผนภาพลิสซาจูและคำนวณกำลังไฟฟ้าจริงเปรียบเทียบกับค่ากำลังไฟฟ้าปรากฏในโปรแกรม MATLAB/Python
    \item ปฏิบัติการวัดสเปกตรัมการเปล่งแสง OES ของพลาสมาเจ็ทที่ป้อนก๊าซอาร์กอนและฮีเลียมที่สภาวะความชื้นต่างกัน
    \item ทดสอบโมเดล 1D-CNN เพื่อทำนายความเข้มข้นของอนุมูลไฮดรอกซิลจากข้อมูลสเปกตรัมดิบแบบเรียลไทม์
\end{enumerate}
\end{planbox}

\begin{planbox}{สื่อการเรียนการสอน}
\begin{enumerate}
    \item แหล่งกำเนิดพลาสมา DBD แผ่นระนาบและหัวพ่นพลาสมาเจ็ทความถี่วิทยุ/พัลส์นาโนวินาที
    \item เครื่องสเปกโตรมิเตอร์ใยแก้วนำแสงความละเอียดสูง (ช่วงความยาวคลื่น 200--900 nm)
    \item สไลด์บรรยายอิเล็กทรอนิกส์และแบบจำลอง 3 มิติแสดงโครงสร้างทางกายภาพของรีแอคเตอร์พลาสมา
    \item โค้ดตัวอย่างภาษาไพทอนสำหรับการสกัดเส้นสเปกตรัมและการคำนวณ Boltzmann Plot
\end{enumerate}
\end{planbox}

\begin{planbox}{การวัดและประเมินผล}
\begin{enumerate}
    \item ประเมินทักษะการปฏิบัติการต่อวงจรและวิเคราะห์ผลลัพธ์ลิสซาจู ร้อยละ 20
    \item ประเมินรายงานผลการวินิจฉัยสเปกตรัม OES และการคำนวณอุณหภูมิการหมุน ร้อยละ 20
    \item ประเมินการทำแบบฝึกหัดทบทวนท้ายบท ร้อยละ 10
    \item ประเมินจากการสอบข้อเขียนประจำภาคการศึกษา ร้อยละ 50
\end{enumerate}
\end{planbox}

\clearpage

% =========================================================================
% เนื้อหาบทเรียน (เริ่มต้นหน้า 3 อย่างเป็นระเบียบ)
% =========================================================================

\section{สถาปัตยกรรมแหล่งกำเนิดพลาสมาเย็นที่ความดันบรรยากาศ}

การนำพลาสมามาประยุกต์ใช้ในงานเกษตรกรรมและสิ่งแวดล้อม จำเป็นต้องใช้ระบบพลาสมาที่ทำงานภายใต้ความดันบรรยากาศปกติ (Atmospheric Pressure Non-Thermal Plasma หรือ APNTP) เพื่อขจัดข้อจำกัดเรื่องห้องสุญญากาศที่มีราคาสูงและไม่เหมาะกับการประมวลผลผลผลิตทางการเกษตรในปริมาณมาก เทคโนโลยีแหล่งกำเนิดพลาสมาที่ได้รับการพัฒนาและใช้งานอย่างแพร่หลายในปัจจุบันประกอบด้วย 4 รูปแบบหลัก

\subsection{การคายประจุแบบแผงกั้นไดอิเล็กทริก}

การคายประจุแบบแผงกั้นไดอิเล็กทริก (Dielectric Barrier Discharge หรือ DBD) ถูกคิดค้นขึ้นครั้งแรกโดย Ernst Werner von Siemens ในปี ค.ศ. 1857 เพื่อการผลิตก๊าซโอโซน โครงสร้างพื้นฐานของระบบ DBD ประกอบด้วยขั้วไฟฟ้าโลหะอย่างน้อย 2 ขั้ว โดยมีชั้นของวัสดุฉนวนไดอิเล็กทริก (Dielectric layer) เช่น แก้วควอตซ์ (Quartz), อลูมินาเซรามิก (\ce{Al2O3}), หรือแก้วไพเร็กซ์ (Pyrex) ขวางกั้นอยู่ระหว่างช่องว่างการคายประจุ (Discharge gap) อย่างน้อยหนึ่งด้าน

หน้าที่สำคัญของแผ่นไดอิเล็กทริกคือการจำกัดปริมาณประจุไฟฟ้าที่สะสมบนผิวฉนวน เมื่อเกิดสตรีมเมอร์ขึ้น ประจุไฟฟ้าจะวิ่งไปสะสมที่ผิวของไดอิเล็กทริก ทำให้เกิดสนามไฟฟ้าย้อนกลับ (Counter-electric field, $E_{ind}$) ที่มีทิศทางตรงข้ามกับสนามไฟฟ้าภายนอก สนามไฟฟ้าย้อนกลับนี้จะลดทอนสนามไฟฟ้ารวมในช่องว่างให้ต่ำกว่าเกณฑ์การพังทลายทางไฟฟ้าอย่างรวดเร็วภายในเวลาไม่กี่นาโนวินาที ป้องกันไม่ให้การคายประจุเปลี่ยนผ่านไปสู่การเกิดอาร์กความร้อนสูง (Thermal Arc)

โครงสร้างทางวิศวกรรมของรีแอคเตอร์ DBD ในงานเกษตรแบ่งออกเป็น 2 ประเภทหลัก
\begin{enumerate}
    \item \textbf{แบบระนาบคู่ขนาน (Planar DBD)} เหมาะสำหรับการปรับสภาพผิวเมล็ดพันธุ์พืช (Seed surface functionalization) แบบสายพานลำเลียง โดยเมล็ดพืชจะไหลผ่านช่องว่างระหว่างแผ่นขั้วไฟฟ้าที่มีการคายประจุสม่ำเสมอ
    \item \textbf{แบบท่อร่วมแกน (Coaxial / Cylindrical DBD)} ประกอบด้วยขั้วไฟฟ้าทรงกระบอกซ้อนกัน เหมาะสำหรับการผลิตน้ำกระตุ้นด้วยพลาสมา (PAW) และการบำบัดก๊าซมลพิษ โดยก๊าซหรือของเหลวจะไหลผ่านตามแนวแกนท่อ
\end{enumerate}

\subsection{หัวพ่นพลาสมาเจ็ทความดันบรรยากาศ}

หัวพ่นพลาสมาเจ็ท (Atmospheric Pressure Plasma Jet หรือ APPJ) เป็นแหล่งกำเนิดพลาสมาที่สามารถนำลำพลาสมาเคลื่อนที่ออกจากขั้วไฟฟ้าไปยังเป้าหมายภายนอก (Effluent zone) ได้อย่างอิสระ ลำพลาสมาที่เห็นพุ่งออกมาด้วยตาเปล่านั้น แท้จริงแล้วไม่ได้เป็นลำก๊าซต่อเนื่อง แต่เป็นกลุ่มก้อนพลาสมาขนาดเล็กมากที่เรียกว่า \textbf{กระสุนพลาสมา (Plasma Bullets)} ซึ่งเคลื่อนที่ด้วยความเร็วสูงยิ่งยวดในระดับ $10^4 - 10^5\text{ m/s}$ ภายใต้อิทธิพลของคลื่นการแตกตัวเป็นไอออนแบบนำวิถี (Guided Ionization Waves)

ข้อได้เปรียบที่โดดเด่นของ APPJ คือความยืดหยุ่นในการนำไปใช้งานเฉพาะจุด (Localized treatment) เช่น การยับยั้งเชื้อราบนแผลเปิดของผลไม้สด การรักษาโรคพืชบนใบ หรือการบำบัดเนื้อเยื่อชีวภาพ โดยที่อุณหภูมิของลำพลาสมาที่สัมผัสผิวชิ้นงานคงที่อยู่ในช่วง $25 - 40^\circ\text{C}$ ทำให้เนื้อเยื่อมีชีวิตไม่เกิดความเสียหายจากความร้อน

\begin{figure}[H]
\centering
\begin{tikzpicture}[scale=0.88, every node/.style={transform shape}]
    % Planar DBD
    \draw[thick, fill=blue!10] (0, 2.4) rectangle (5, 2.7);
    \node at (2.5, 2.9) {\textbf{ขั้วไฟฟ้าบน (High Voltage)}};
    
    \draw[thick, fill=orange!30] (0, 2.0) rectangle (5, 2.4);
    \node at (2.5, 2.2) {\small ชั้นฉนวนไดอิเล็กทริก (Alumina / Quartz)};
    
    % Discharge Gap
    \fill[cyan!15] (0, 0.7) rectangle (5, 2.0);
    \draw[dashed, cyan!80!blue, thick] (0.5, 0.8) -- (0.5, 1.9);
    \draw[dashed, cyan!80!blue, thick] (1.5, 0.8) -- (1.5, 1.9);
    \draw[dashed, cyan!80!blue, thick] (2.5, 0.8) -- (2.5, 1.9);
    \draw[dashed, cyan!80!blue, thick] (3.5, 0.8) -- (3.5, 1.9);
    \draw[dashed, cyan!80!blue, thick] (4.5, 0.8) -- (4.5, 1.9);
    \node at (2.5, 1.35) {\textbf{\color{blue!80!black}ช่องว่างคายประจุพลาสมา (Discharge Gap)}};
    
    \draw[thick, fill=gray!30] (0, 0.3) rectangle (5, 0.7);
    \node at (2.5, 0.5) {\textbf{ขั้วไฟฟ้าล่าง (Ground)}};
    
    \draw[thick, ->] (-0.8, 1.35) -- (-0.1, 1.35);
    \node at (-1.5, 1.35) {\small ทางเข้าก๊าซ};
    \draw[thick, ->] (5.1, 1.35) -- (5.8, 1.35);
    \node at (6.6, 1.35) {\small ก๊าซที่ถูกกระตุ้น};
    
    \node at (2.5, -0.2) {\textbf{(ก) โครงสร้างรีแอคเตอร์แบบแผงกั้นไดอิเล็กทริก (Planar DBD)}};
\end{tikzpicture}
\caption{สถาปัตยกรรมโครงสร้างพื้นฐานของระบบการคายประจุแบบแผงกั้นไดอิเล็กทริก}
\label{fig:dbd_schematic}
\end{figure}

\section{การวินิจฉัยทางไฟฟ้าและแผนภาพลิสซาจู}

เนื่องจากระบบพลาสมาเย็นเป็นโหลดทางไฟฟ้าแบบไม่เชิงเส้น (Non-linear load) ที่มีการเกิดและดับของไมโครดิสชาร์จในระดับนาโนวินาที การวัดกำลังไฟฟ้าด้วยเครื่องวัดโวลต์-แอมป์มิเตอร์ทั่วไปจึงเกิดความคลาดเคลื่อนสูงมาก วิศวกรรมพลาสมาจึงอาศัย \textbf{แผนภาพลิสซาจู (Lissajous Figure หรือ $Q-V$ plot)} เป็นวิธีมาตรฐานในการคำนวณกำลังงานที่ถ่ายโอนลงสู่พลาสมาอย่างแท้จริง

\subsection{วงจรสมมูลของแมนเลย์}

ตามแบบจำลองวงจรสมมูลของแมนเลย์ (Manley's Equivalent Circuit) ระบบ DBD ประกอบด้วยตัวเก็บประจุ 2 ตัวต่ออนุกรมกัน
\begin{enumerate}
    \item \textbf{ความจุไฟฟ้าของแผ่นไดอิเล็กทริก ($C_d$)} ซึ่งเป็นค่าคงที่ตามโครงสร้างเรขาคณิตและชนิดของวัสดุฉนวน
    \item \textbf{ความจุไฟฟ้าของช่องว่างก๊าซ ($C_g$)} เมื่อก๊าซยังไม่แตกตัว ค่าความจุนี้จะมีค่าตามค่าคงที่ไดอิเล็กทริกของอากาศ ($\varepsilon_r \approx 1$) แต่เมื่อเกิดการพังทลายทางไฟฟ้า ช่องว่างจะกลายเป็นตัวนำไฟฟ้าที่มีความต้านทานพลาสมา $R_p$ ต่อขนานอยู่
\end{enumerate}

เมื่อทำการต่อตัวเก็บประจุสำหรับตรวจวัดประจุ ($C_m$ โดยทั่วไปเลือกใช้ค่า $10 - 100\text{ nF}$) อนุกรมเข้ากับขั้วกราวด์ของรีแอคเตอร์ ประจุสะสมรวม $Q(t)$ สามารถคำนวณได้จากแรงดันตกคร่อมตัวเก็บประจุตรวจวัด $V_m(t)$ ดังสมการ
\begin{equation}
    Q(t) = C_m \cdot V_m(t)
\end{equation}

\subsection{การคำนวณกำลังไฟฟ้าจริงจากพื้นที่รูปสี่เหลี่ยมด้านขนาน}

เมื่อนำสัญญาณแรงดันไฟฟ้ากระตุ้น $V(t)$ มาพล็อตลงบนแกน $X$ และสัญญาณประจุ $Q(t)$ ลงบนแกน $Y$ ของออสซิลโลสโคป จะได้กราฟรูปสี่เหลี่ยมด้านขนานปิด (Lissajous parallelogram) ซึ่งพื้นที่ปิดภายในเส้นโค้งนี้แสดงถึงพลังงานไฟฟ้าที่สูญเสียไปในการสร้างพลาสมาต่อหนึ่งรอบคลื่น ($E_{cycle}$)
\begin{equation}
    E_{cycle} = \oint Q(t) \, dV(t)
    \label{eq:lissajous_energy}
\end{equation}
ดังนั้น กำลังไฟฟ้าเฉลี่ยที่ถ่ายโอนลงสู่พลาสมา ($P_{plasma}$) จึงมีค่าเท่ากับผลคูณของพลังงานต่อรอบคลื่นกับความถี่ใช้งาน ($f$)
\begin{equation}
    P_{plasma} = f \cdot E_{cycle} = f \oint Q \, dV
    \label{eq:lissajous_power}
\end{equation}

ความชันของแต่ละด้านในรูปสี่เหลี่ยมด้านขนานลิสซาจูยังบ่งบอกคุณสมบัติทางไฟฟ้าที่สำคัญ โดยด้านที่พลาสมาดับจะมีความชันเท่ากับความจุไฟฟ้ารวมของระบบ ($C_{cell} = \frac{C_d C_g}{C_d + C_g}$) ส่วนด้านที่พลาสมาติดจะมีความชันเท่ากับความจุไฟฟ้าของแผ่นไดอิเล็กทริก ($C_d$)

\section{การวินิจฉัยด้วยสเปกโทรสโกปีการเปล่งแสง}

การวินิจฉัยด้วยสเปกโทรสโกปีการเปล่งแสง (Optical Emission Spectroscopy หรือ OES) เป็นเทคนิคการวินิจฉัยพลาสมาแบบไม่รบกวน (Non-invasive diagnostics) ที่มีประสิทธิภาพสูงสุด โดยอาศัยหลักการตรวจวัดโฟตอนของแสงที่ปลดปล่อยออกมาเมื่ออะตอม โมเลกุล หรือไอออนที่ถูกกระตุ้นเกิดการเปลี่ยนระดับพลังงานจากสถานะกระตุ้นกลับสู่สถานะที่ต่ำกว่า

\subsection{การระบุสปีชีส์อนุภาค RONS ในพลาสมา}

ในสเปกตรัม OES ของพลาสมาความดันบรรยากาศที่ใช้ในงานเกษตร แถบสเปกตรัมที่สำคัญประกอบด้วย
\begin{enumerate}
    \item \textbf{อนุมูลไฮดรอกซิล ($\cdot\ce{OH}$)} เปล่งแสงที่ความยาวคลื่น $309\text{ nm}$ จากการเปลี่ยนผ่านระดับพลังงาน $A^2\Sigma^+ \rightarrow X^2\Pi$ เป็นตัวบ่งชี้ถึงความสามารถในการออกซิเดชันสูงของระบบ
    \item \textbf{ระบบโพสิทีฟที่สองของไนโตรเจน (\ce{N2} Second Positive System, SPS)} เปล่งแสงในช่วงความยาวคลื่น $315 - 380\text{ nm}$ จากการเปลี่ยนสถานะ $C^3\Pi_u \rightarrow B^3\Pi_g$
    \item \textbf{ระบบเนกาทีฟแรกของไอออนโมเลกุลไนโตรเจน (\ce{N2^+} First Negative System, FNS)} เปล่งแสงที่ความยาวคลื่น $391.4\text{ nm}$ จากการเปลี่ยนสถานะ $B^2\Sigma_u^+ \rightarrow X^2\Sigma_g^+$ บ่งชี้ถึงการมีอยู่ของอิเล็กตรอนพลังงานสูงที่มากกว่า $18.7\text{ eV}$
    \item \textbf{ออกซิเจนอะตอมิก (\ce{O})} เปล่งแสงที่เส้นสเปกตรัมเดี่ยวความยาวคลื่น $777.4\text{ nm}$ ($3p^5P \rightarrow 3s^5S^\circ$) และ $844.6\text{ nm}$ ($3p^3P \rightarrow 3s^3S^\circ$)
\end{enumerate}

\subsection{การหาอุณหภูมิการหมุนและอุณหภูมิของก๊าซ}

ในพลาสมาเย็นที่ความดันบรรยากาศ อัตราการถ่ายโอนพลังงานระหว่างการเคลื่อนที่แบบหมุนของโมเลกุลและการชนแบบยืดหยุ่นของโมเลกุลก๊าซเกิดขึ้นอย่างรวดเร็วมาก (ระดับ $10^{-9}\text{ s}$) ส่งผลให้อุณหภูมิการหมุนของโมเลกุล (Rotational Temperature, $T_{rot}$) มีค่าอยู่ในสมดุลกับอุณหภูมิจลน์ของก๊าซ (Gas Temperature, $T_g$) นั่นคือ $T_g \approx T_{rot}$

การคำนวณหา $T_{rot}$ สามารถทำได้โดยใช้วิธี \textbf{โบลต์ซมันน์พล็อต (Boltzmann Plot Method)} จากเส้นสเปกตรัมการหมุนของโมเลกุล $\cdot\ce{OH}$ หรือ \ce{N2} ตามสมการ
\begin{equation}
    \ln\left( \frac{I_{J'J''}}{S_{J'J''} \nu^4} \right) = -\frac{E_{rot}(J')}{k_B T_{rot}} + \text{const}
    \label{eq:boltzmann_plot}
\end{equation}
โดยที่ $I_{J'J''}$ คือความเข้มแสงของเส้นสเปกตรัม, $S_{J'J''}$ คือแฟกเตอร์เฮินล์-ลอนดอน (Hönl-London factor), $\nu$ คือความถี่ของแสง, $E_{rot}(J')$ คือพลังงานการหมุนของสถานะบน, และ $k_B$ คือค่าคงที่โบลต์ซมันน์ เมื่อพล็อตความสัมพันธ์ระหว่างเทอมทางซ้ายมือกับพลังงาน $E_{rot}$ ความชันของกราฟเส้นตรงจะเท่ากับ $-\frac{1}{k_B T_{rot}}$ ทำให้สามารถหาอุณหภูมิของก๊าซได้อย่างแม่นยำโดยไม่ต้องสัมผัสกับชิ้นงาน

\begin{aibox}{การประยุกต์ใช้ 1D-CNN และ Transformer ในการแยกส่วนสเปกตรัม OES}
ในทางปฏิบัติ สเปกตรัม OES ที่ตรวจวัดได้จากหัวพ่นพลาสมาจริงมักประสบปัญหาสัญญาณรบกวน (Noise), การเลื่อนของเส้นฐาน (Baseline Drift), และการซ้อนทับกันอย่างหนาแน่นของแถบการหมุน (Spectral Overlapping) การใช้วิธีวิเคราะห์แบบเดิมเช่น การปรับเทียบด้วยเกาส์เซียนหรือลอเรนซ์เซียนต้องใช้เวลาประมวลผลนานหลายสิบวินาทีต่อหนึ่งสเปกตรัม จึงไม่สามารถตอบสนองต่อการควบคุมแบบทันทีได้

\textbf{สถาปัตยกรรม 1D-CNN และ Spectral Transformer:} ได้รับการออกแบบขึ้นเพื่อแก้ปัญหานี้โดยเฉพาะ
\begin{enumerate}
    \item \textbf{1D Convolutional Neural Network (1D-CNN):} ทำหน้าที่สกัดคุณลักษณะเฉพาะของยอดคลื่น (Feature Extraction) และกำจัดสัญญาณรบกวนความถี่สูง โดยใช้ชั้นคอนโวลูชันขนาดฟิลเตอร์ต่างๆ (Multi-scale Kernels)
    \item \textbf{Multi-Head Self-Attention:} โมดูล Transformer ทำหน้าที่จับความสัมพันธ์ระยะไกลข้ามช่วงความยาวคลื่น เพื่อประเมินสภาวะการกระจายตัวของพลังงานการสั่นและการหมุน
\end{enumerate}
\begin{equation}
    \widehat{T}_{rot}, \widehat{T}_{vib}, [\cdot\ce{OH}] = \mathcal{M}_{deep}(\mathbf{S}_{raw}; \mathbf{\Theta}^*)
\end{equation}
โมเดลนี้สามารถสกัดค่าอุณหภูมิ $T_g$ และความเข้มข้นของอนุมูลอิสระสำคัญได้ภายในเวลาเพียง \textbf{0.8 มิลลิวินาที} ด้วยความแม่นยำสูงกว่าร้อยละ 98.4 เมื่อเทียบกับโปรแกรมจำลองมาตรฐานสากล (Specair) ทำให้เปิดทางไปสู่การควบคุมพลาสมาเจ็ทแบบปรับตัวได้ (Adaptive Plasma Jet Controller) ซึ่งจะปรับแรงดันไฟฟ้าและความถี่อัตโนมัติเพื่อป้องกันไม่ให้พืชได้รับความร้อนเกินเกณฑ์ความปลอดภัย
\end{aibox}

\begin{examplebox}{ตัวอย่างการคำนวณหากำลังไฟฟ้าของระบบ DBD จากแผนภาพลิสซาจู}
\textbf{โจทย์} ระบบกำเนิดพลาสมา DBD แผ่นระนาบทำงานที่ความถี่กระตุ้นไซนูซอยด์ $f = 15\text{ kHz}$ เมื่อต่อตัวเก็บประจุตรวจวัดประจุ $C_m = 47\text{ nF}$ ออสซิลโลสโคปบันทึกแผนภาพลิสซาจูได้เป็นรูปสี่เหลี่ยมด้านขนานที่มีพิกัดจุดยอด 4 จุด บนระนาบ $(V, V_m)$ ดังนี้
\begin{itemize}
    \item จุด $A: (-5.0\text{ kV}, -20\text{ V})$
    \item จุด $B: (-2.0\text{ kV}, +15\text{ V})$
    \item จุด $C: (+5.0\text{ kV}, +20\text{ V})$
    \item จุด $D: (+2.0\text{ kV}, -15\text{ V})$
\end{itemize}
จงคำนวณหา (ก) พลังงานที่สูญเสียในการคายประจุต่อหนึ่งรอบคลื่น ($E_{cycle}$) และ (ข) กำลังไฟฟ้าเฉลี่ยที่ถ่ายโอนลงสู่พลาสมา ($P_{plasma}$)

\textbf{การแก้ปัญหา}
\begin{enumerate}
    \item แปลงค่าแรงดัน $V_m$ เป็นประจุไฟฟ้าสะสม $Q = C_m \cdot V_m$
    \begin{align*}
        Q_A &= 47 \times 10^{-9}\text{ F} \times (-20\text{ V}) = -0.94\text{ }\mu\text{C} \\
        Q_B &= 47 \times 10^{-9}\text{ F} \times (+15\text{ V}) = +0.705\text{ }\mu\text{C} \\
        Q_C &= 47 \times 10^{-9}\text{ F} \times (+20\text{ V}) = +0.94\text{ }\mu\text{C} \\
        Q_D &= 47 \times 10^{-9}\text{ F} \times (-15\text{ V}) = -0.705\text{ }\mu\text{C}
    \end{align*}
    \item คำนวณพื้นที่ของรูปสี่เหลี่ยมด้านขนานลิสซาจู $ABCD$
    \begin{equation*}
    \begin{aligned}
        \text{Area} &= \frac{1}{2} \big| V_A(Q_B - Q_D) + V_B(Q_C - Q_A) \\
        &\quad + V_C(Q_D - Q_B) + V_D(Q_A - Q_C) \big|
    \end{aligned}
    \end{equation*}
    แทนค่าตัวเลขพิกัด $(V \text{ ในหน่วย V}, Q \text{ ในหน่วย C})$
    \begin{equation*}
    \begin{aligned}
        E_{cycle} &= 2 \times \left[ (V_C - V_A) \times \frac{(Q_B - Q_D)}{2} \right] \times \frac{1}{2} \\
        &= 10000\text{ V} \times 1.41 \times 10^{-6}\text{ C} \times \frac{1}{2} \\
        &= 7.05 \times 10^{-3}\text{ J} = 7.05\text{ mJ}
    \end{aligned}
    \end{equation*}
    \item คำนวณกำลังไฟฟ้าเฉลี่ย $P_{plasma} = f \cdot E_{cycle}$
    \begin{equation*}
        P_{plasma} = 15 \times 10^3\text{ s}^{-1} \times 7.05 \times 10^{-3}\text{ J} = 105.75\text{ W}
    \end{equation*}
\end{enumerate}
\textbf{ความหมายทางกายภาพ} ระบบ DBD นี้ใช้กำลังไฟฟ้าเฉลี่ยจริงในการแตกตัวไอออนและสังเคราะห์สารเคมีพลาสมาเท่ากับ 105.75 วัตต์ ซึ่งสามารถนำไปใช้คำนวณประสิทธิภาพเชิงพลังงานในการผลิตปุ๋ยไนเตรตหรือการกำจัดเชื้อจุลินทรีย์ต่อไป
\end{examplebox}

\clearpage

% =========================================================================
% คำถามทบทวนท้ายบทที่ 3 (ขึ้นหน้าใหม่อย่างสง่างาม)
% =========================================================================

\section*{คำถามทบทวนท้ายบทที่ 3}
\addcontentsline{toc}{section}{คำถามทบทวนท้ายบทที่ 3}

\begin{enumerate}
    \item จงอธิบายกลไกทางฟิสิกส์ที่ทำให้แผ่นไดอิเล็กทริกในระบบ DBD สามารถยับยั้งการเปลี่ยนผ่านจากไมโครดิสชาร์จไปสู่อาร์กความร้อนได้ และหากเปลี่ยนวัสดุจากแก้วควอตซ์เป็นอลูมินาเซรามิกจะส่งผลต่อความจุไฟฟ้าของระบบอย่างไร
    \item ในการวิเคราะห์แผนภาพลิสซาจู ($Q-V$ plot) ความชันของเส้นกราฟช่วงที่พลาสมาดับและช่วงที่พลาสมาติดสัมพันธ์กับค่าความจุไฟฟ้าของชิ้นส่วนใดในวงจรสมมูลของแมนเลย์ จงแสดงความสัมพันธ์ทางคณิตศาสตร์
    \item เหตุใดการวินิจฉัยสเปกตรัมการเปล่งแสง (OES) จึงเป็นที่นิยมอย่างยิ่งในการศึกษาพลาสมาเย็นสำหรับการเกษตร เมื่อเทียบกับการใช้โพรบแบบสัมผัส เช่น โพรบแลงเมียร์ (Langmuir Probe)
    \item เส้นสเปกตรัมที่ความยาวคลื่น 309 นาโนเมตร และ 391.4 นาโนเมตร เกิดจากการเปลี่ยนผ่านระดับพลังงานของสปีชีส์เคมีใดในพลาสมา และเส้นสเปกตรัมทั้งสองนี้บ่งชี้ถึงคุณสมบัติทางฟิสิกส์ใดของพลาสมา
    \item จงเปรียบเทียบข้อดีและข้อจำกัดของการใช้อัลกอริทึม Deep Learning (1D-CNN/Transformer) เทียบกับวิธีดั้งเดิมในการประมวลผลสัญญาณ OES สำหรับการควบคุมสภาวะการทำงานของหัวพ่นพลาสมาเจ็ทในโรงงานบรรจุผลผลิตทางการเกษตร
\end{enumerate}

\clearpage
